% !PRINTED 555
\weberpage{555}
%<<CONT
liegt, so ist die Endlichkeit der Anzahl der Classen nachgewiesen.

Lassen wir, wie bisher, $\omega_1, \omega_2, \ldots, \omega_n$ eine Minimalbasis von $\Omega$ bedeuten, so werden alle ganzen Zahlen des Körpers aus

\weq{1}{\omega = \omega_1 x_1 + \omega_2 x_2 + \cdots + \omega_n x_n}

erhalten, wenn wir den $x_1, x_2, \ldots, x_n$ ganze rationale Zahlwerthe ertheilen. Wenn wir eine positive ganze Zahl $k$ annehmen und festsetzen, dass keine der Zahlen $x_i$ aus dem Intervall $\pm k$ heraustreten soll, so wird der \gesperrt{absolute Werth} von $\omega$ nicht über der Grenze

\[
(|\omega_1| + |\omega_2| + \cdots + |\omega_n|) k = rk
\]

liegen, wenn unter $|\omega_1|, |\omega_2|, \ldots$ die absoluten Werthe (Einleitung, Bd. I, S. 18) der (reellen oder complexen) Grössen $\omega_i$ verstanden sind, und $r$ die Summe $|\omega_1| + |\omega_2| + \cdots$ bedeutet. Bilden wir den Ausdruck (1) für die $n$ conjugirten Körper, und nehmen das Product, so erhalten wir, wenn wir mit $R$ eine positive reelle Zahl bezeichnen, die über dem Product der $n$ Werthe $r$ liegt,

\weq{2}{N_a(\omega) < R k^n.}

Diese Zahl $R$ ist nur von der Natur des Körpers $\Omega$, nicht aber von $k$ abhängig.

Jetzt sei $\mu$ irgend ein ganzes Functional in $\Omega$, und $N_a(\mu)$ seine absolute Norm. Wenn wir die ganze Zahl $k$ so bestimmen, dass

\weq{3}{k^n \leqq N_a(\mu) < (k+1)^n,}

und wenn wir ferner in (1) den Zahlen $x_i$ die Werthe $0, 1, 2, \ldots, k$ ertheilen, so ist die Anzahl der verschiedenen Werthe, die aus (1) hervorgehen, $(k+1)^n$, also grösser als $N_a(\mu)$. Nach \S.\,148, 3. ist aber die Zahl der nach dem Modul $\mu$ incongruenten Zahlen gleich $N_a(\mu)$, und folglich müssen unter den so bestimmten Zahlen $\omega$ mindestens zwei verschiedene nach dem Modul $\mu$ congruente Zahlen vorkommen. Ist also $\omega' \equiv \omega'' (\mathrm{mod}\ \mu)$, so wird die Differenz

\weq{4}{\alpha = \omega' - \omega'' = (x_1' - x_1'') \omega_1 + \cdots + (x_n' - x_n'') \omega_n}

durch $\mu$ theilbar sein, und zugleich sind die ganzen Zahlen

\[
x_1' - x_1'' = a_1, \ldots, x_n' - x_n'' = a_n
\]
%CONT>>
